\documentclass[12pt,a4paper]{article}

\usepackage[T1]{fontenc}
\usepackage[utf8]{inputenc}
\usepackage{mathptmx}
\usepackage[a4paper,margin=2.5cm]{geometry}
\usepackage{setspace}
\usepackage{microtype}
\usepackage[french]{babel}
\usepackage{csquotes}
\usepackage{graphicx}
\usepackage{xcolor}
\usepackage{pifont}
\usepackage{booktabs}
\usepackage{tabularx}
\usepackage{array}
\usepackage{float}
\usepackage{tikz}
\usetikzlibrary{arrows.meta,babel}
\usepackage{etoolbox}
\usepackage[nobottomtitles*]{titlesec}
\renewcommand{\bottomtitlespace}{5\baselineskip}
\usepackage{caption}
\makeatletter
\setlength{\@fptop}{4pt}
\setlength{\@fpsep}{8pt plus 1fil}
\setlength{\@fpbot}{0pt plus 1fil}
\makeatother
\usepackage{xurl}
\usepackage[hidelinks]{hyperref}

\onehalfspacing
\setlength{\parindent}{1.25em}
\setlength{\parskip}{0pt}

\captionsetup{font=small,labelfont=bf,skip=6pt}
\titleformat{\section}{\large\bfseries}{\thesection.}{0.6em}{}
\titleformat{\subsection}{\normalsize\bfseries}{\thesubsection}{0.6em}{}
\titlespacing*{\section}{0pt}{1.4em}{0.7em}
\titlespacing*{\subsection}{0pt}{1.1em}{0.4em}

\hypersetup{
  pdftitle={Alerte de proximité avant la crue à Kinshasa},
  pdfauthor={}
}

\AtBeginEnvironment{tikzpicture}{\shorthandoff{:;!?}}
\tikzset{
  box/.style={draw=black, fill=white, align=center, inner sep=3.5pt, font=\scriptsize},
  databox/.style={draw=black, fill=white, align=left, inner sep=3.5pt, font=\scriptsize},
  layer/.style={draw=black!70, fill=black!4, inner sep=0pt},
  arr/.style={draw=black, line width=0.7pt, -{Stealth[length=2.2mm]}},
  linklbl/.style={draw=black, fill=white, font=\scriptsize\bfseries, inner sep=2.5pt, align=center}
}

\begin{document}

\begin{titlepage}
\hypersetup{pageanchor=false}
\thispagestyle{empty}
\centering
\vspace*{2mm}
{\bfseries RÉPUBLIQUE DÉMOCRATIQUE DU CONGO\par}
\vspace{0.35em}
{\bfseries MINISTÈRE DE L'ENSEIGNEMENT SUPÉRIEUR ET UNIVERSITAIRE\par}
\vspace{0.7em}
{\Large\bfseries UNIVERSITÉ DE KINSHASA\par}
\vspace{0.35em}
\includegraphics[height=2.8cm]{figures/logo-unikin.png}\\[0.35em]
{\Large\bfseries Faculté des Sciences et Technologies\par}
\vspace{0.25em}
{\large Mention Mathématiques, Statistique et Informatique\par}
\vspace{0.2em}
{\bfseries Licence Master et Doctorat\par}
\vspace{1.4cm}

\fcolorbox{cyan!20}{blue!35}{%
\begin{minipage}{0.96\textwidth}
\centering
\vspace{0.85cm}
{\LARGE\sffamily ALERTE DE PROXIMITÉ AVANT LA CRUE\\[0.35em]
UNE ARCHITECTURE IOT POUR LES QUARTIERS\\[0.25em]
RIVERAINS DE LA NDJILI ET DU POOL MALEBO\\[0.25em]
À KINSHASA\par}
\vspace{0.85cm}
\end{minipage}%
}

\vspace{0.6em}
{\Huge\textcolor{blue}{\ding{45}}\par}
\vspace{0.35em}
{\LARGE\sffamily\bfseries KALEMBU MASIRIKA Constantin\par}

\vspace{1.1cm}
\begin{flushright}
\begin{tabular}{lcl}
\multicolumn{3}{l}{\textbf{\textsf{Travail pratique de récupération}}}\\
\multicolumn{3}{l}{\textbf{\textsf{Internet des objets (IoT)}}}\\[0.55em]
\textbf{\textsf{Promotion}} & \textbf{\textsf{:}} & \textbf{\textsf{Master 1 --- Data Science et IA}}\\[0.45em]
\textbf{\textsf{Enseignant}} & \textbf{\textsf{:}} & \textbf{\textsf{Michel Tshimanga}}\\
\end{tabular}
\end{flushright}

\vfill
{\Large\sffamily\bfseries Année académique 2025 -- 2026\par}
\vspace{0.8cm}
\end{titlepage}

\hypersetup{pageanchor=true}
\tableofcontents
\clearpage

\section{Introduction}

Le thème du travail est large. L'Internet des objets, mis devant les problèmes de développement de la République démocratique du Congo, peut accueillir à peu près n'importe quel capteur. J'ai choisi un cas où l'eau est déjà passée, et où l'on voit précisément ce qui a manqué.

Dans la nuit du 4 au 5 avril 2025, des pluies torrentielles font déborder la rivière Ndjili. Treize des vingt-quatre communes de Kinshasa sont touchées. Le bilan repris par la Fédération internationale de la Croix-Rouge fait état de 165~morts, de 28~blessés et de plus de 7\,000~personnes déplacées, avec des hauteurs d'eau allant jusqu'à un mètre et demi dans des quartiers périphériques d'une ville de plus de dix-sept millions d'habitants~\cite{ifrc2025}. L'année précédente, en janvier 2024, des inondations qualifiées de record avaient déjà concerné au moins six communes et plus de 79\,000~personnes, pour des dommages directs estimés à 96~millions de dollars~\cite{priur2026}.

La METTELSAT, l'Agence nationale de météorologie et de télédétection par satellite, avait annoncé des pluies supérieures à la normale plusieurs semaines avant avril 2025~\cite{ifrc2025}. L'avertissement saisonnier existait donc. Ce qui n'existait pas, c'est une mesure du niveau, proche des parcelles, capable de dire à un chef de quartier, le matin même, qu'un repère connu vient d'être dépassé.

C'est cet écart que la proposition cherche à réduire. Les sections suivent l'ordre demandé. D'abord le problème, tel qu'il se pose à Kinshasa, et non sous la forme d'une ville intelligente imaginaire. Ensuite le pilote, avec un schéma d'architecture original. Puis un tableau comparatif de trois technologies de communication, lu pour un choix et non pour une liste. Ensuite deux risques de déploiement, chacun avec une parade, et deux limites que ces parades ne lèvent pas. La conclusion dit ce que je retiens, et ce que je refuse de faire dire à ce schéma.

\section{Un problème déjà mesuré, mal annoncé}
\label{sec:probleme}

Kinshasa est bâtie entre les collines et le Pool Malebo. Les termes de référence de l'étude de drainage lancée dans le Projet de résilience aux inondations urbaines (PRIUR) rappellent que la ville est drainée par plusieurs rivières. Les plus importantes sont la Ndjili, dont le bassin versant est d'environ 2\,000~km\textsuperscript{2}, et la N'sele, d'environ 6\,000~km\textsuperscript{2}. Des cours plus courts comptent autant dans les quartiers~: Gombe, Funa, Kalamu, Makelele, et d'autres~\cite{priur2026}. Ces lits sont l'exutoire. Quand ils sont encombrés, la pluie reste dans les parcelles.

Le 11 avril 2025, devant le Conseil des ministres, le ministre des Infrastructures et Travaux publics décrit un réseau devenu inopérant. Les rejets réduisent le lit de la Ndjili, de la Funa, de la Matete, de la Kalamu, de la Makelele et de la Gombe. Les berges et les zones non aedificandi, c'est-à-dire les zones où l'on ne doit pas bâtir, sont occupées. L'ouest de la ville, autour de la baie de Ngaliema, court le même risque. Les remèdes qu'il cite ne sont pas des capteurs~: des travaux sur les rivières, le dragage du fleuve, la gestion des déchets, et l'arrêt des constructions dans le lit majeur~\cite{conseil2025}.

Il faut distinguer deux eaux, parce qu'un seul type de mesure les confondrait. La première est rapide. Une nuit de pluie sur le bassin, et la Ndjili monte en quelques heures, comme les 4 et 5 avril 2025. La seconde est lente. En janvier 2024, l'OCHA indiquait que le fleuve Congo avait atteint, selon la Régie des voies fluviales, son plus haut niveau depuis plus de soixante ans~\cite{ocha2024}. Le Pool monte, l'évacuation des rivières urbaines se bloque, et des quartiers prennent l'eau sans qu'un orage local suffise à l'expliquer. Kinsuka et la baie de Ngaliema relèvent plutôt de cette seconde eau. Masina et Kingabwa relèvent des deux.

Les suites ne s'arrêtent pas au niveau dans la maison. En avril 2025, la Croix-Rouge signale des coupures d'électricité et, à Kinshasa, une interruption de l'eau de la REGIDESO pendant cinq jours~\cite{ifrc2025}. L'eau stagnante vient ensuite. Le système dont je parle n'a pas à prétendre traiter une épidémie. Il doit arriver avant que les biens, et les personnes, soient déjà dans l'eau.

Une partie du travail de fond est déjà engagée, et il serait malhonnête de l'ignorer. La METTELSAT produit les prévisions. Le projet CREWS, mené par la Banque mondiale et l'Organisation météorologique mondiale avec la METTELSAT comme partenaire opérationnel, appuie ces services hydrométéorologiques et la préparation d'outils de gestion du risque, y compris des termes de référence liés au drainage de Kinshasa~\cite{crews2025}. Le PRIUR, financé par un crédit de l'Association internationale de développement de 200~millions de dollars, vise d'abord Kinshasa et Kalemie. Sur ce montant, 150~millions vont aux infrastructures résilientes et 30~millions au renforcement institutionnel de la gestion des risques de catastrophe~\cite{priur2026}. Des relevés de rue ont même déjà cartographié l'exposition de Matete, de N'Djili et de Lemba~\cite{bm2020}.

Tout cela répond à d'autres questions~: quel temps fera la saison, quels ouvrages reconstruire, quelles parcelles sont exposées. Aucun de ces dispositifs ne dit, à l'échelle d'un marché de Masina, que le niveau vient de passer la dalle que les vendeurs connaissent. Le trou est là. Il est local, il porte sur les heures qui viennent, et il doit joindre des gens qui n'ouvriront pas une application.

\section{Ce que ferait le pilote}
\label{sec:pilote}

Je limite le pilote à trois stations de niveau. Chacune porte aussi un pluviomètre.

\begin{itemize}
  \item Une sur la Ndjili, à l'entrée de la zone urbaine de Masina, pour voir la crue arriver.
  \item Une plus en aval, du côté de Kingabwa, là où le lit urbain est le plus contraint.
  \item Une au bord du Pool, vers Kinsuka, pour la montée du fleuve et le blocage de l'évacuation.
\end{itemize}

La figure~\ref{fig:sites} rappelle que ces trois points ne mesurent pas la même eau. Les positions sont de principe. Les emplacements réels se choisiraient avec la commune, la Régie des voies fluviales et les riverains, sur un ouvrage déjà gardé~: un pont, un bâtiment communal, un centre de santé. Pas un boîtier seul au milieu d'une berge.

\begin{figure}[ht]
\centering
\begin{tikzpicture}[x=1cm,y=1cm,font=\footnotesize]
  \fill[black!5] (0,2.05) rectangle (14.5,6.55);
  \fill[black!9] (0,0) rectangle (14.5,2.05);
  \draw[black, line width=1.5pt, line cap=round]
    (12.4,6.55) .. controls (11.95,4.7) and (10.75,3.2) .. (9.55,2.05);
  \node[anchor=south east, font=\footnotesize\bfseries] at (12.2,6.2) {Ndjili};
  \node[databox] (mas) at (9.35,5.25) {\textbf{Station Masina}\\amont urbain\\niveau + pluie\\pas 15~min / 2~min};
  \draw (mas.east) -- (11.65,5.25);
  \node[databox] (kg) at (7.05,3.45) {\textbf{Station Kingabwa}\\aval urbain\\niveau + pluie\\pas 15~min / 2~min};
  \draw (kg.east) -- (10.35,3.35);
  \node[anchor=north west, align=left, font=\scriptsize] at (0.25,6.35) {\textbf{Bassin + ville}\\Pluie locale (auget)\\vs onde de crue.};
  \node[anchor=north west, align=left, font=\scriptsize] at (0.25,4.25) {\textbf{Crue rapide}\\La Ndjili monte\\en quelques heures.};
  \node[databox] (ks) at (2.15,0.95) {\textbf{Station Kinsuka}\\bord du Pool\\niveau + pluie\\blocage d'évacuation};
  \node[anchor=west, font=\footnotesize\bfseries] at (4.55,1.25) {Pool Malebo};
  \node[anchor=west, align=left, font=\scriptsize] at (8.0,0.95) {\textbf{Crue lente}\\Le fleuve monte et\\gêne la sortie des rivières.};
  \node[anchor=south west, font=\scriptsize] at (0.2,2.15) {Zone rivières urbaines};
  \node[anchor=north west, font=\scriptsize] at (0.2,1.95) {Zone fleuve / Pool};
\end{tikzpicture}
\caption{Deux eaux, trois stations. Schéma de principe, pas un levé. Chaque station mesure le niveau (hors d'eau) et la pluie.}
\label{fig:sites}
\end{figure}

Chaque station mesure sans toucher l'eau. Un capteur à ultrason ou à radar, fixé au-dessus du lit, évite l'envasement et les déchets qui arrachent une sonde immergée. Le pas ordinaire est de quinze minutes. Si le niveau varie vite, il passe à deux minutes. Le pluviomètre sert à distinguer une averse locale d'une onde qui descend la rivière. Le nœud n'envoie qu'une cote, une lame d'eau et l'état de sa batterie.

Les seuils ne sortent pas d'un modèle que personne n'a calé sur ces sections. Ils partent de repères que les habitants utilisent déjà~: une dalle de marché, une marche de pont, un mur de parcelle. On les relève une fois, on les associe à la cote du capteur, et on les nomme dans le message. Quatre états suffisent. Ils sont dans le tableau~\ref{tab:etats}.

\begin{table}[ht]
\centering
\small
\renewcommand{\arraystretch}{1.28}
\begin{tabularx}{\textwidth}{@{}>{\raggedright\arraybackslash\bfseries}p{2.15cm}>{\raggedright\arraybackslash}X>{\raggedright\arraybackslash}X>{\raggedright\arraybackslash}X@{}}
\toprule
État & Repère & Ce qui part & Qui décide \\
\midrule
Veille & Sous le repère connu du site & Rien vers le public. La cote est seulement archivée. & La station, seule. \\
\addlinespace
Vigilance & Le niveau approche du repère & SMS aux relais de ce site seulement. & Automatique, dans la passerelle. \\
\addlinespace
Alerte & Le repère est dépassé & SMS aux relais. Sirène locale si le lien vers le centre est coupé. & Automatique, dans la passerelle. \\
\addlinespace
Crue majeure & Le repère historique est dépassé & Demande d'évacuer les parcelles basses de ce site. & Une personne responsable confirme. \\
\bottomrule
\end{tabularx}
\caption{Quatre états par station. Seul le dernier attend un visa humain. Les trois premiers ne donnent pas l'ordre de quitter le quartier.}
\label{tab:etats}
\end{table}

Cette séparation est le cœur du dispositif. La vigilance et l'alerte demandent de mettre les biens en hauteur, pas de vider un quartier. Elles peuvent partir sans attendre un bureau. La crue majeure, elle, dit aux gens de partir. Un automate qui lance cet ordre au premier capteur déréglé fera plus de mal que le silence. Le serveur propose. Une personne désignée confirme~: la protection civile provinciale, ou le bourgmestre, selon ce que la ville aura écrit avant la saison. Pas le soir de la crue.

Le message est court. Il dit le lieu, le repère, le geste, et vers qui se tourner. Il est prévu en français et en lingala.

\begin{quote}
\textbf{Exemple de message, état \enquote{alerte}, station de Masina.}

\textbf{Français.} Ndjili, Masina. Niveau au-dessus du repère orange. L'eau peut entrer dans les parcelles basses. Mettez les biens en hauteur. Confirmez auprès du chef de quartier.

\vspace{0.45em}
\textbf{Lingala, base à reprendre sur place.} Mayi ya Ndjili na Masina ekoli. Repère orange eleki. Bomema biloko na likolo. Tuna na chef ya quartier.
\end{quote}

\vspace{0.4em}
La phrase lingala n'est pas un texte définitif. \enquote{Mayi ekoli} dit que l'eau monte, \enquote{bomema biloko na likolo} demande de monter les biens, \enquote{tuna na chef ya quartier} renvoie au chef de quartier. Elle doit être reprise avec les relais avant tout envoi. Un mot de trop, ou un ton de rumeur, et le deuxième message ne sera plus cru. Le français du SMS a le même devoir de brièveté~: un paragraphe ne s'affiche pas sur un téléphone simple, et il ne se crie pas dans une allée.

Le relais humain fait partie du système, au même titre que la radio. Le chef de quartier, le responsable du marché et l'infirmier titulaire atteignent déjà des parcelles sans smartphone et sans forfait. Le SMS les vise eux d'abord. Une page peut exister pour la mairie, afin de voir les trois cotes. Elle n'est pas le canal qui prévient le quartier.

\section{Architecture}
\label{sec:archi}

L'architecture tient en quatre couches. Toutes sont dans la ville. La figure~\ref{fig:archi} détaille le parcours de l'information~: capteurs, protocole, décision locale, serveur, SMS et relais humains.

\begin{figure}[p]
\centering
\begin{tikzpicture}[x=1cm,y=1cm,font=\scriptsize]
  \fill[black!5] (0.05,15.55) rectangle (15.55,19.55);
  \fill[black!8] (0.05,11.15) rectangle (15.55,15.35);
  \fill[black!5] (0.05,6.85) rectangle (15.55,10.95);
  \fill[black!8] (0.05,0.15) rectangle (15.55,6.65);
  \draw[black!55] (0.05,15.45) -- (15.55,15.45);
  \draw[black!55] (0.05,11.05) -- (15.55,11.05);
  \draw[black!55] (0.05,6.75) -- (15.55,6.75);

  \node[font=\small\bfseries] at (7.8,19.95) {Architecture IoT du pilote --- alerte de proximité à Kinshasa};

  \node[anchor=west, font=\footnotesize\bfseries] at (0.2,19.25) {1. Perception (terrain)};
  \node[databox, text width=4.35cm] (n1) at (2.7,17.55) {%
    \textbf{Station Masina}\\
    Ndjili, amont urbain\\
    Capteur niveau (ultrason/radar)\\
    Pluviomètre à auget\\
    Batterie + boîtier fermé\\
    Pas~: 15~min (2~min si $\Delta$ rapide)};
  \node[databox, text width=4.35cm] (n2) at (7.8,17.55) {%
    \textbf{Station Kingabwa}\\
    Ndjili, aval urbain\\
    Capteur niveau (ultrason/radar)\\
    Pluviomètre à auget\\
    Batterie + boîtier fermé\\
    Pas~: 15~min (2~min si $\Delta$ rapide)};
  \node[databox, text width=4.35cm] (n3) at (12.9,17.55) {%
    \textbf{Station Kinsuka}\\
    Bord du Pool Malebo\\
    Capteur niveau (ultrason/radar)\\
    Pluviomètre à auget\\
    Batterie + boîtier fermé\\
    Pas~: 15~min (2~min si $\Delta$ rapide)};
  \node[anchor=west, align=left] at (0.25,15.85) {Charge utile~: cote (m), lame d'eau (mm), état batterie, horodatage. Pose sur un ouvrage déjà gardé.};

  \draw[arr] (7.8,15.55) -- (7.8,15.05);
  \node[linklbl] at (11.2,15.3) {LoRaWAN $\approx$\,868~MHz\\portée utile 1--3~km (à mesurer)};

  \node[anchor=west, font=\footnotesize\bfseries] at (0.2,14.95) {2. Collecte et décision locale};
  \node[databox, text width=9.4cm] (gw) at (5.4,13.05) {%
    \textbf{Passerelle --- bâtiment communal gardé}\\
    Réception LoRaWAN des 3 stations\\
    Archivage local si le centre est injoignable\\
    Comparaison cote $\leftrightarrow$ repère du quartier\\
    États~: veille / vigilance / alerte / crue majeure\\
    Vigilance et alerte~: automatiques pour ce site};
  \node[databox, dashed, text width=4.4cm] (sir) at (12.85,13.05) {%
    \textbf{Sirène / haut-parleur}\\
    Chemin de secours\\
    Part si le lien vers le\\
    serveur est coupé\\
    Décidée avant la saison};
  \node[anchor=west, align=left] at (0.25,11.45) {Règle locale~: une station seule peut demander \enquote{biens en hauteur}. Elle n'ordonne pas seule une évacuation.};

  \draw[arr] (7.8,11.15) -- (7.8,10.65);
  \node[linklbl] at (11.4,10.9) {Lien urbain + file d'attente\\données originales à Kinshasa};

  \node[anchor=west, font=\footnotesize\bfseries] at (0.2,10.55) {3. Traitement et supervision (dans la ville)};
  \node[databox, text width=7.0cm] (srv) at (4.1,8.55) {%
    \textbf{Serveur à Kinshasa}\\
    Hôtel de ville / établissement public\\
    Historique des cotes\\
    Journal~: seuil, heure, destinataires\\
    Proposition d'évacuation (crue majeure)\\
    Original des données conservé ici};
  \node[databox, text width=6.7cm] (hum) at (12.0,8.55) {%
    \textbf{Visa humain}\\
    Protection civile ou bourgmestre\\
    Confirme seulement la crue majeure\\
    Exige 2 stations d'accord\\
    ou 1 station + appel de contrôle\\
    Règle écrite avant la saison};
  \node[anchor=west, align=left] at (0.25,7.15) {Tableau de bord mairie (consultation). L'alerte ne dépend jamais d'un lien international.};

  \draw[arr] (7.8,6.85) -- (7.8,6.35);
  \node[linklbl] at (11.25,6.6) {SMS (FR + lingala)\\2 opérateurs (ex.~Vodacom / Airtel)};

  \node[anchor=west, font=\footnotesize\bfseries] at (0.2,6.05) {4. Diffusion et action humaine};
  \node[databox, text width=4.5cm] at (2.7,4.15) {%
    \textbf{Chefs de quartier}\\
    Relais vers les parcelles\\
    Confirment ou démentent\\
    sur place};
  \node[databox, text width=4.5cm] at (7.8,4.15) {%
    \textbf{Marchés}\\
    Biens et vivres\\
    mis en hauteur\\
    Avant l'entrée d'eau};
  \node[databox, text width=4.5cm] at (12.9,4.15) {%
    \textbf{Centres de santé}\\
    Salles basses et stock\\
    Mis à l'abri\\
    Infirmier titulaire};
  \node[databox, text width=14.8cm] at (7.8,1.85) {%
    \textbf{Contenu type du SMS (état alerte)} --- lieu + repère + geste + contact.\\
    FR~: \textit{Ndjili, Masina. Niveau au-dessus du repère orange. Mettez les biens en hauteur.}\\
    Lingala (base)~: \textit{Mayi ya Ndjili na Masina ekoli. Bomema biloko na likolo. Tuna na chef ya quartier.}};
  \node[anchor=west, align=left] at (0.25,0.55) {Liste courte de destinataires nommés. Pas d'envoi de masse à toute la commune le soir d'une crue.};
\end{tikzpicture}
\caption{Schéma original de l'architecture IoT proposée. Quatre couches, données échangées, protocoles et règles de décision. Le cadre tireté est le chemin de secours local.}
\label{fig:archi}
\end{figure}

\paragraph{Perception.}
Le nœud est un objet contraint~: un microcontrôleur, une radio longue portée, le capteur de niveau, un pluviomètre à auget, une batterie tenue pour plusieurs jours sans soleil. Un petit panneau n'est ajouté que si le toit qui le porte est déjà surveillé. Le boîtier est fermé, sans écran, sans câble qui pend. Il n'a pas l'air d'un objet à revendre.

\paragraph{Collecte et décision locale.}
Une passerelle reçoit les trois stations de sa zone. Elle est dans un bureau qui a déjà un gardien. Elle conserve les mesures quand le lien vers le centre est coupé, et c'est elle qui compare la cote au repère. Le seuil d'alerte peut donc sonner sur place, par une sirène ou un haut-parleur, sans attendre un serveur. Ce choix s'écrit avant la saison des pluies, avec la commune. Il ne s'improvise pas le jour où l'eau monte.

\paragraph{Traitement.}
Le serveur est une machine ordinaire, physiquement à Kinshasa~: à l'hôtel de ville, ou dans une salle d'un établissement public qui accepte d'en répondre, l'université par exemple. Il montre les cotes, garde l'historique, et tient un journal simple. Quelle cote, quel seuil, qui a confirmé le message de crue majeure, à quelle heure le SMS est parti. L'original reste dans la ville. Une copie peut exister ailleurs, pour ne pas tout perdre dans un incendie ou un vol de machine. Le jour où le lien international est instable, l'alerte ne doit pas dépendre de cette copie.

\paragraph{Diffusion.}
Le SMS part d'une puce, avec une seconde puce chez un autre opérateur. On retient, pour le quartier visé, les deux réseaux qui passent encore sous la pluie. On le vérifie en saison, pas sur une brochure. Vodacom, Airtel et Orange sont les candidats à comparer sur place. La liste des destinataires reste courte~: les relais nommés, pas \enquote{toute la commune}. Un envoi de masse, le soir d'une crue, fabrique la saturation que l'on voulait éviter.

\paragraph{Action.}
Les relais font ensuite ce qu'ils font déjà. Ils parlent au quartier. Le système leur donne une heure et un repère, pas un discours.

Deux règles empêchent un nœud isolé de faire peur trop loin. L'alerte \enquote{biens en hauteur} peut venir d'une seule station, parce qu'elle ne concerne que le voisinage de cette station. La demande de confirmation pour une évacuation exige soit deux stations d'accord, Masina et Kingabwa par exemple, soit une station et un appel de vérification. Un boîtier arraché qui envoie une cote absurde ne suffit pas à vider un quartier.

\clearpage
\section{Tableau comparatif de trois technologies de communication}
\label{sec:tech}

L'énoncé demande un tableau comparatif de trois technologies de communication. Je compare LoRaWAN, le NB-IoT et le GSM via SMS~\cite{lora,uit2018}, parce que le pilote a deux liaisons distinctes~: remonter une cote depuis le lit, puis joindre un téléphone déjà présent dans le quartier. Le tableau~\ref{tab:tech} les confronte critère par critère.

\begin{table}[H]
\centering
\scriptsize
\renewcommand{\arraystretch}{1.08}
\begin{tabularx}{\textwidth}{@{}>{\raggedright\arraybackslash\bfseries}p{2.85cm}>{\raggedright\arraybackslash}X>{\raggedright\arraybackslash}X>{\raggedright\arraybackslash}X@{}}
\toprule
Critère & LoRaWAN & NB-IoT & GSM / SMS \\
\midrule
Nature & LPWAN privé, topologie en étoile (nœuds $\rightarrow$ passerelle). & LPWAN cellulaire (3GPP), sur réseau d'opérateur. & Réseau mobile classique~; le SMS sert de canal d'alerte. \\
\addlinespace
Standard & LoRa Alliance~\cite{lora}. & 3GPP (NB-IoT). & GSM / 3GPP, service SMS. \\
\addlinespace
Spectre & Bande sub-GHz des dispositifs à courte portée (souvent $\approx$\,868~MHz, Région~1 UIT). Confirmation ARPTC requise~\cite{uit2018,eleos}. & Bandes licenciées LTE de l'opérateur. & Bandes licenciées (800, 900, 1\,800~MHz, etc.)~\cite{arptc2025}. \\
\addlinespace
Portée à Kinshasa & Environ 1--3~km en bâti dense~; plus si l'antenne est dégagée. À mesurer. & Celle du LTE, seulement si le service objet est activé. & Large en ville, inégale à l'intérieur et sous l'orage. \\
\addlinespace
Débit / charge utile & Faible (octets à quelques centaines d'octets)~: cote, pluie, batterie. & Faible à moyen~: adapté aux objets. & SMS $\approx$\,160 caractères~: message d'alerte bref. \\
\addlinespace
Consommation & Très faible~: le nœud dort presque tout le temps. & Faible, mais modem + SIM par station. & Élevée pour un modem permanent au bord de l'eau. \\
\addlinespace
Latence & Secondes à minutes. Acceptable pour une crue. & Secondes. & Quelques secondes à une minute. \\
\addlinespace
Infrastructure & Passerelle(s) installée(s) et entretenue(s) localement. & Couverture NB-IoT + abonnement par objet. & Réseau déjà déployé~; 1--2 puces côté serveur. \\
\addlinespace
Coût d'exploitation & Installation et maintenance de la passerelle~; pas d'abonnement par station. & Abonnement mensuel par capteur, toute l'année. & Forfait + SMS pour 1--2 puces~; trop cher par station. \\
\addlinespace
Sous orage & Lien local, indépendant du réseau mobile. & Même fragilité que le réseau cellulaire. & Peut saturer ou couper au moment utile~\cite{ifrc2025}. \\
\addlinespace
Disponibilité RDC & Possible après vérification ARPTC. & Offre objet commerciale non démontrée à Kinshasa. & Déjà présente (Vodacom, Airtel, Orange, etc.). \\
\addlinespace
Adaptation au cas & Bonne pour remonter les mesures~; mauvaise pour joindre un relais. & Possible sur le papier~; dépend d'un service non vérifié. & Mauvaise pour le capteur~; bonne pour joindre les relais. \\
\addlinespace
Décision & \textbf{Retenu}~: capteur $\rightarrow$ passerelle. & \textbf{Non retenu} pour ce pilote. & \textbf{Retenu}~: passerelle $\rightarrow$ relais (SMS). \\
\bottomrule
\end{tabularx}
\caption{Tableau comparatif de trois technologies de communication pour l'alerte de proximité à Kinshasa.}
\label{tab:tech}
\end{table}

\subsection{Lecture du tableau}

Le tableau ne sert pas à élire une technologie unique. Il sépare deux rôles. LoRaWAN et le NB-IoT se ressemblent sur le papier~: faible débit, faible consommation, longue portée. Ils divergent sur l'essentiel en RDC. LoRaWAN demande une passerelle à installer, mais ne dépend pas d'une offre commerciale d'opérateur. Le NB-IoT dépend précisément de cette offre. Or une bande LTE attribuée n'est pas un service objet ouvert~\cite{arptc2025}. Retenir le NB-IoT aujourd'hui, ce serait bâtir le pilote sur une hypothèse non démontrée à Kinshasa, plus une carte SIM dans chaque boîtier au bord de l'eau.

Le GSM via SMS n'est pas un concurrent de LoRaWAN sur le même tronçon. Il est mauvais comme lien de capteur~: consommation, vol de puce, coût multiplié, avenir incertain du 2G. Il est bon comme dernier kilomètre humain. Une ou deux puces, dans un bâtiment fermé, joignent des téléphones simples. C'est le seul des trois déjà dans la poche des chefs de quartier.

La résilience sous orage tranche le mélange. En avril 2025, l'électricité et l'eau de Kinshasa ont cédé ensemble~\cite{ifrc2025}. Le réseau mobile peut céder aussi. Le lien LoRaWAN station--passerelle reste local. Le SMS doit être doublé (deux opérateurs) et complété par une sirène locale. Le NB-IoT hériterait de la même fragilité cellulaire.

La bande LoRaWAN visée est celle des dispositifs à courte portée en dessous du gigahertz, souvent autour de 868~MHz~\cite{uit2018}. Une compilation privée indique, pour la RDC, la plage 868--868,6~MHz à 25~mW~\cite{eleos}. Ce n'est pas une décision de l'ARPTC. Avant tout achat, la bande et la puissance se vérifient auprès de l'Autorité. Si elle impose une autre canalisation, on change la fréquence des modules, pas le choix architectural du tableau~\ref{tab:tech}.

Le choix final est donc mixte. LoRaWAN entre le lit et la maison communale. Le SMS ensuite, vers les personnes. Le NB-IoT pourra remplacer l'un ou l'autre le jour où un opérateur le proposera à Kinshasa avec un tarif d'objet et une couverture vérifiée sous la pluie. Ce jour n'est pas un fait, aujourd'hui.

\section{Deux risques au moment du déploiement}
\label{sec:risques}

Un risque, ici, est un événement qui peut faire échouer le pilote ou le rendre dangereux, et contre lequel une parade existe dans la façon de le poser. Les limites de la section suivante sont autre chose~: elles restent vraies même si le pilote fonctionne.

\subsection{Le matériel disparaît, ou il se met à mentir}

Un panneau visible, une antenne, un boîtier au bord d'une avenue~: cela se démonte. Le risque n'est pas seulement la perte du matériel. Un nœud ouvert peut servir à émettre une fausse cote. Une fausse alerte, deux fois de suite, et les chefs de quartier cessent de relayer. Le système devient alors pire que l'absence de système, parce qu'il a usé la seule chose qui le faisait tenir~: la confiance.

La parade consiste à réduire l'exposition, pas à ajouter de la tôle. Les stations sont sur des ouvrages déjà gardés. Le panneau solaire n'est posé que sur un toit surveillé. Le boîtier n'a ni écran ni valeur de revente évidente. Un petit stock de nœuds de rechange est prévu dès le départ~: le vol est un événement attendu, pas une surprise qui arrête la saison. Côté radio, l'entrée sur le réseau suit le mécanisme de session et de clés prévu par LoRaWAN, plutôt qu'un émetteur ouvert que n'importe qui peut imiter~\cite{lora}. La puce SMS, elle, reste dans le bâtiment, avec un code, pas dans le boîtier au-dessus de l'eau.

Surtout, une station seule ne peut pas provoquer le message d'évacuation. Elle peut au plus demander, pour son propre voisinage, de mettre les biens en hauteur. Ce message cite le repère et l'heure. Un relais qui voit le ciel clair et la rivière basse peut le contredire tout de suite, et cette contradiction doit remonter dans le journal. Mieux vaut un SMS démenti par le chef de quartier qu'une consigne que personne n'ose contredire.

\subsection{La panne au moment où la pluie tombe}

C'est le risque qui compte. Le pilote est inutile s'il meurt avec le courant, et nuisible s'il parle dans le vide pendant que le réseau mobile sature. Avril 2025 a montré que les grands réseaux de Kinshasa peuvent céder ensemble~: électricité, et eau de la REGIDESO pendant cinq jours~\cite{ifrc2025}. Rien n'autorise à penser que le réseau mobile, lui, tiendra le soir d'une crue de la Ndjili.

La passerelle tient plusieurs jours sur sa batterie. La comparaison entre la cote et le repère se fait dans cette passerelle, pas sur une machine lointaine. Si le centre ne répond plus, la sirène convenue à l'avance reste possible. Le serveur est dans la ville~: une coupure du câble international ne bloque pas l'envoi. Le SMS part par deux opérateurs, parce qu'un réseau saturé ne doit pas être le chemin unique. On teste cet envoi un jour de forte pluie, avant le pic de saison, et l'on note lequel des deux passe encore. Enfin, la liste des destinataires reste courte. Vouloir prévenir tout Kinshasa par SMS le soir d'une crue, c'est organiser soi-même l'encombrement.

\section{Deux limites que le schéma ne lève pas}
\label{sec:limites}

Une limite n'est pas un incident. C'est une borne. La parade du risque ne la supprime pas. En nommer deux oblige à ne pas vendre le pilote pour autre chose que ce qu'il est.

\subsection{L'alerte ne fait pas baisser l'eau}

Même si chaque SMS part à l'heure, la Ndjili débordera tant que son lit reste encombré, tant que les caniveaux sont bouchés, et tant que l'on construit dans les zones où l'on ne doit pas bâtir. Le ministre des Infrastructures le disait en avril 2025 sans parler d'objets connectés~: les travaux, le dragage, les déchets, la discipline sur les berges~\cite{conseil2025}. Le PRIUR oriente 150~millions de dollars vers les infrastructures, sur un financement de 200~millions~\cite{priur2026}. Poser un capteur à côté de ce chantier et laisser entendre que la ville est couverte serait une faute.

Le pilote fait gagner du temps. Il ne remplace ni un curage, ni un exutoire, ni une parcelle hors d'eau. Si l'alerte sert d'alibi pour retarder ces travaux, il vaut mieux ne pas l'installer. Je la défends seulement comme un filet, en attendant des ouvrages qui, eux, retirent l'eau. Le filet n'est pas l'ouvrage.

\subsection{Trois stations voient peu, et personne n'est payé pour les tenir}

Masina, Kingabwa et Kinsuka ne décrivent pas treize communes. Une averse sur la Funa ou la Kalamu peut noyer un quartier pendant que la Ndjili, en amont, est encore calme. Inversement, Kinsuka peut monter avec le fleuve sans que Masina soit déjà touchée. Le message doit donc dire de quel site il parle, et se taire sur le reste. Étendre le réseau à toutes les rivières de la ville est une autre décision, avec d'autres passerelles et d'autres relais. Ce n'est pas un réglage.

L'autre borne est plus ordinaire, et elle suffit à tuer le projet. Les capteurs dérivent. Les augets se bouchent avec la poussière et les feuilles. Les batteries se fatiguent. Les numéros des chefs de quartier changent. Quelqu'un doit marcher jusqu'à la station, quelqu'un doit payer les deux puces, quelqu'un doit avoir le droit de valider la crue majeure. Si ces trois fonctions n'ont pas de budget de fonctionnement après la cérémonie d'installation, le pilote devient un boîtier muet en une saison des pluies. Je ne connais pas de correctif technique à cette limite. C'est une question de contrat~: qui emploie le technicien en novembre, et sur quelle ligne de la ville.

\clearpage
\section{Conclusion}

Je retiens l'alerte de proximité comme un complément étroit. La prévision de saison existe déjà, elle est portée par la METTELSAT, et des partenaires l'appuient. L'argent des ouvrages est engagé dans le PRIUR. Le trou utile est plus bas, et il est concret~: la cote d'un lit, un repère que le quartier connaît déjà, un message bref en français et en lingala vers des relais qui font déjà circuler les nouvelles, et une sirène qui ne dépend ni du courant du réseau urbain ni d'une machine hors du pays.

Je ne retiens pas l'idée qu'un capteur pourrait remplacer l'autorité qui décide d'évacuer. Je ne retiens pas non plus qu'une application remplacerait le chef de quartier, le marché ou le centre de santé. Ces trois-là sont déjà l'infrastructure. Le pilote ne fait que leur donner une heure et une cote. Je ne retiens pas, enfin, un déploiement sur les vingt-quatre communes tant que trois stations n'ont pas passé une saison des pluies entière, vol, boue et pannes compris.

Le test n'est pas le schéma. Au cœur de la saison, le test est de savoir qui change la batterie, qui paie les deux puces, et qui signe le message quand le repère historique est dépassé. Sans ces trois noms écrits quelque part, l'architecture reste un dessin. Kinshasa n'a pas besoin d'un dessin de plus. Elle a besoin d'un avertissement qui part encore la deuxième année.

\clearpage
\phantomsection
\addcontentsline{toc}{section}{Références}
\begin{thebibliography}{99}
\bibitem{ifrc2025}
Fédération internationale des sociétés de la Croix-Rouge et du Croissant-Rouge.
\emph{DREF Operation MDRCD046, Democratic Republic of the Congo, Floods}.
2025.
\href{https://reliefweb.int/attachments/d10b84da-d4c9-45a3-aaa3-0594b96b4c72/MDRCD046do.pdf}{reliefweb.int, opération DREF MDRCD046}.

\bibitem{ocha2024}
Bureau de la coordination des affaires humanitaires.
\emph{Flash Update n\textsuperscript{o}~1~: de graves inondations affectent 18~provinces}.
République démocratique du Congo, 24 janvier 2024.
\href{https://reliefweb.int/attachments/54eda70e-28fb-4f52-915e-dbef7e99789c/OCHA\%20DRC\%20Flash\%20Update1\%20Inondations\%202023-2024.pdf}{reliefweb.int, flash update du 24 janvier 2024}.

\bibitem{conseil2025}
Présidence de la République démocratique du Congo.
\emph{Compte rendu de la trente-neuvième réunion ordinaire du Conseil des ministres du 11 avril 2025}.
Communication du ministre des Infrastructures et Travaux publics sur le débordement de la Ndjili.
\href{https://presidence.cd/uploads/files/COMPTE\%20RENDU\%20DE\%20LA\%20TRENTE-NEUVIEME\%20REUNION\%20ORDINAIRE\%20DU\%20CONSEIL\%20DES\%20MINISTRES\%20DU\%2011\%20AVRIL\%202025.pdf}{presidence.cd, compte rendu du 11 avril 2025}.

\bibitem{priur2026}
Projet de résilience aux inondations urbaines en République démocratique du Congo.
\emph{Termes de référence pour l'étude du réseau global de drainage de la ville de Kinshasa}.
Financement IDA, Banque mondiale. Document de consultation diffusé en 2026.
\href{https://www.mediacongo.net/dpics/filesmanager/jobaof/2026/02-fevrier/01/ci_260205_02_04.pdf}{mediacongo.net, termes de référence PRIUR}.

\bibitem{crews2025}
CREWS, Banque mondiale et Organisation météorologique mondiale.
\emph{Democratic Republic of the Congo. Project status report, January--December 2025}.
Partenaire opérationnel~: METTELSAT.
\href{https://crews-initiative.org/wp-content/uploads/2026/09/CREWS-Democratic-Republic-of-Congo-Project-Status-Report-Jan\%E2\%80\%93Dec-2025.pdf}{crews-initiative.org, rapport 2025}.

\bibitem{bm2020}
Banque mondiale.
\emph{Flood exposure and vulnerability mapping using street level image automatic acquisition}.
Relevés de rue à Matete, N'Djili et Lemba, 5--11 février 2020.
\href{https://documents1.worldbank.org/curated/en/099648304032428971/pdf/IDU1e7b0c081137c5142051b4f017e8c4d6964b9.pdf}{documents.worldbank.org, relevés de rue 2020}.

\bibitem{uit2018}
Union internationale des télécommunications.
\emph{Rapport UIT-R SM.2423~: aspects techniques et opérationnels des réseaux LPWAN pour les communications de type machine et l'Internet des objets}.
2018.
\href{https://www.itu.int/dms_pub/itu-r/opb/rep/R-REP-SM.2423-2018-PDF-F.pdf}{itu.int, rapport SM.2423}.

\bibitem{lora}
LoRa Alliance.
Spécification LoRaWAN, couche liaison.
\href{https://lora-alliance.org/about-lorawan/}{lora-alliance.org}.

\bibitem{arptc2025}
Autorité de régulation de la poste et des télécommunications du Congo.
\emph{Questionnaire de consultation publique sur les réseaux 5G et le spectre}.
Mars 2025. Le document décrit les bandes mobiles en discussion. Il ne décrit pas une offre NB-IoT commerciale.
\href{https://arptc.gouv.cd/app/uploads/2025/03/Questionnaire_Consultation-publique_RDC_V7_010325.pdf}{arptc.gouv.cd, consultation de mars 2025}.

\bibitem{eleos}
Eleos Compliance.
Fiche de synthèse sur l'homologation et des plages de fréquences en République démocratique du Congo.
Compilation privée, à ne pas confondre avec une décision de l'ARPTC.
\href{https://www.eleoscompliance.com/fr/type-approval/congo-republique-d\%C3\%A9mocratique-du-congo}{eleoscompliance.com, fiche RDC}.
\end{thebibliography}

\end{document}
